\subsection{The Floating-Point Representation Architecture (\texttt{float})}

Python's \texttt{float} type represents real-number \emph{approximations}, not exact real numbers. An \texttt{int} can represent arbitrarily large whole numbers exactly. A \texttt{float} uses a fixed-size binary format, so many decimal fractions can only be approximated.

\subsubsection{Architectural Realities: Mapping Python Floats to C Doubles via the IEEE 754 Double-Precision Binary Standard}

In CPython, the ordinary Python \texttt{float} is implemented as the platform C \texttt{double}. On modern machines, that means IEEE~754 double-precision binary floating point: 64 bits total, approximately 53 bits of significand precision, fixed exponent range.

Python exposes limits through \texttt{sys.float\_info}:

\begin{itemize}
    \item \texttt{radix = 2}: binary representation;
    \item \texttt{mant\_dig = 53}: about 53 binary digits of precision;
    \item \texttt{dig = 15}: about 15 reliable decimal digits;
    \item \texttt{max}: largest finite value; overflow yields \texttt{inf}.
\end{itemize}

Unlike \texttt{int}, \texttt{float} does not grow. \texttt{sys.float\_info.max * 2} produces \texttt{inf}, not a larger exact number. The representation is fixed-width at the hardware level.

\subsubsection{Machine Epsilon and Precision Loss: The Pitfalls of Representing Base-10 Fractions in Base-2 Memory Buffers}

The canonical surprise:

\begin{verbatim}
0.1 + 0.2 == 0.3        # False
0.1 + 0.2               # 0.30000000000000004
\end{verbatim}

This is not a Python bug. Decimal fractions such as \texttt{0.1} lack finite binary expansions, just as \texttt{1/3} lacks a finite decimal expansion. The literal \texttt{0.1} is converted to the closest representable IEEE~754 value. \texttt{float.as\_integer\_ratio()} and \texttt{float.hex()} expose the stored binary fraction directly.

Machine epsilon (\texttt{sys.float\_info.epsilon}) marks the smallest increment distinguishable from \texttt{1.0} at that scale. For computed float equality, \texttt{math.isclose()} is appropriate; direct \texttt{==} is often wrong.

\begin{quote}
Use \texttt{float} for approximate real-number computation, not for exact decimal accounting.
\end{quote}

\subsubsection{Special Floating-Point Values: Infinity, Negative Infinity, and NaN}

Floating-point arithmetic includes special values: \texttt{inf}, \texttt{-inf}, and \texttt{nan}. All remain \texttt{float} objects at the Python level.

\begin{verbatim}
float("nan") == float("nan")   # False
math.isnan(float("nan"))       # True
\end{verbatim}

NaN is not equal even to itself. Test with \texttt{math.isnan()}, not \texttt{==}. Use \texttt{math.isinf()} and \texttt{math.isfinite()} for the other special cases.

The practical summary:

\begin{quote}
Python's \texttt{float} is a fixed-precision IEEE~754 binary object mapped from C \texttt{double}. It is fast and useful for approximate computation, but cannot represent all decimal fractions exactly, can overflow to infinity, and includes non-orderable NaN values.
\end{quote}
